\section{Evaluation Design}
\label{sec:evaluation}

This section specifies the planned evaluation. It does not report final \toolname results. All missing results are explicitly marked.

\subsection{Dataset}

The main dataset will be the benchmark used by ``Vulnerability-Affected Versions Identification: How Far Are We?''~\cite{howfar2025}. The local reference analysis states that the benchmark contains 1,128 C/C++ vulnerabilities from nine open-source projects and 132 vulnerability types. Exact project-level table values and patch-type counts must be verified against repaired tables before final use \need{NEEDS TABLE REPAIR}. The dataset is selected because it is directly aligned with our target problem: given CVEs and patch commits, identify affected release versions.

\subsection{Baselines}

We will compare against the baselines organized in the benchmark study:
\begin{itemize}[leftmargin=*]
  \item \textbf{Tracing-based:} VCCFinder~\cite{vccfinder2015}, V-SZZ~\cite{vszz2022}, Lifetime~\cite{lifetime2021}, SEM-SZZ~\cite{semszz2024}, TC-SZZ~\cite{tcszz2024}, and LLM4SZZ~\cite{llm4szz2025}.
  \item \textbf{Matching-based:} ReDeBug~\cite{redebug2012}, VUDDY~\cite{vuddy2017}, MOVERY~\cite{movery2022}, V1SCAN~\cite{v1scan2023}, FIRE~\cite{fire2024}, and VULTURE~\cite{vulture2025}.
\end{itemize}

If official runnable artifacts are unavailable or inconsistent, we will report whether baseline values are taken from the paper table or reproduced locally. No reproduced baseline claim will be made until scripts, data, and environment are verified \need{NEEDS ARTIFACT}.

\subsection{Metrics}

We follow the metric structure extracted from the baseline and related papers:
\begin{itemize}[leftmargin=*]
  \item \textbf{CVE-level Accuracy:} the fraction of CVEs for which the predicted affected-version set exactly equals the ground truth set.
  \item \textbf{CVE-level No-Miss Ratio:} the fraction of CVEs for which all ground-truth affected versions are included in the prediction.
  \item \textbf{Version-level Precision/Recall/F1:} flattened version-level metrics over all CVE-version pairs.
\item \textbf{Resource metrics:} average time, input tokens, output tokens, total tokens, cost, agent calls, candidate boundary events, and execution failure rate for \toolname \need{NEEDS EXPERIMENT}.
  \item \textbf{Manual-workload and lower-bound metrics:} if complete FN validation is infeasible for some cases, we will report conservative lower-bound recall and manual inspection counts, following the evaluation style of developer-log based affected-version work~\cite{tdsc2024}.
\end{itemize}

% \input{tables/evaluation_metrics}  % TODO: create this table file

\subsection{Research Questions}

\textbf{RQ1: Overall effectiveness.} How accurately does \toolname identify affected versions compared with tracing-based and matching-based baselines? We will report CVE-level Accuracy, CVE-level NMR, and version-level precision/recall/F1. Result table: \need{NEEDS EXPERIMENT}.

\textbf{RQ2: Robustness across patch and repository conditions.} How does performance vary across add-only, delete-only, mixed patches, single-file and multi-file patches, and single-branch versus multi-branch settings? Subset definitions will follow the benchmark where available. Result table: \need{NEEDS EXPERIMENT}.

\textbf{RQ3: Contribution of each \toolname stage.} What is the effect of patch-family filtering, root-cause evidence extraction, boundary-event judgment, and Git-DAG state propagation? Planned ablations include disabling Step1 filtering, replacing Step2 VET with raw touched-code evidence, replacing repository-aware release projection with simpler tag ordering, and replacing boundary-event judgment with deterministic matching. Result table: \need{NEEDS EXPERIMENT}.

\textbf{RQ4: Cost and failure analysis.} What is the runtime, token cost, and failure rate of \toolname? Which cases fail because of patch noise, root-cause extraction error, boundary-event judgment error, release projection error, timeout, or missing artifacts? Result table: \need{NEEDS EXPERIMENT}.

\textbf{RQ5: Agent behavior and evidence use.} When \toolname invokes an agent, which evidence sources and tool calls drive semantic judgments, and how do candidate-set size, token usage, and cost correlate with success? This RQ is inspired by recent agentic SZZ analyses that report tool calls, tokens, cost, and behavior distributions~\cite{agentszz2026,howwhyagents2026}. Result table: \need{NEEDS EXPERIMENT}.

\subsection{Pilot Status}

The local SystemDesign directory contains current pilot artifacts over a 32-case subset. These artifacts are useful for debugging the paper pipeline and metric implementation, but they are not the final ICSE results and should not be compared directly against full 1,128-CVE paper baselines. Any pilot table copied into the paper must be labeled as a pilot and kept separate from the main evaluation \need{NEEDS EXPERIMENT}.
